%!TEX root = ../summaries.tex
% ******************************* Thesis Appendix A: Glossaries ****************************

\chapter{Research Summaries}

\section*{Influence Maximization}

Maximizing the Spread of Influence through a Social Network \cite{kempe2003}:
This paper formally introduced the problem of Influence Maximization and is first attempt to show that greedy approximation method outperforms centrality and other heuristic based methods in influence mining on two fundamental influence models LT and IC.
	
Cost-effective Outbreak Detection in Networks \cite{leskovec2007b}: generalizes the problem of influence maximization by using sub-modularity property, while proposing CELF, an efficient algorithm 700 times faster than simple greedy technique.
	
On minimizing budget and time in influence propagation over social networks \cite{goyal2012}: proposed MINTSS to reduce the size of initial seed; the second contribution, MINTIME minimizes the time taken to activate minimum required active nodes.
	
Scalable Influence Maximization for Prevalent Viral Marketing in Large-Scale Social Networks \cite{chen2010b}: a rather important study that proposed a tuneable algorithm, MIA to control influence spread versus runtime, thus making the influence maximization possible on large networks.
	
A Data-Based Approach to Social Influence Maximization \cite{goyal2011}: introduced a tempo credit distribution model, which looks into past action log to estimate influence spread. This study also shows that models, which do not learn probabilities from propagation traces suffer from high error rate.
	
Efficient Influence Maximization in Social Networks \cite{chen2009}: this study challenges the conclusion from \cite{kempe2003} that greedy approximation outperforms heuristics, by proposing a heuristic called degree discount that runs hundreds of times faster and show almost similar results. The experiments thoroughly run on CELF, Greedy, Mixed Greedy and Degree Discount algorithms on IC model and show the efficiency of proposed heuristic.
	
Scalable Influence Maximization in Social Networks under the Linear Threshold Model \cite{chen2010b}: takes an unusual approach to solve influence maximization problem by computing DACs (directed acyclic graphs) in linear time on LT model. 
	
Time-Critical Influence Maximization in Social Networks with Time-Delayed Diffusion Process \cite{chen2012}: a time constrained method, which allows influence maximization under time critical settings. This study introduced two heuristics. The first one saves influence computation in a tree and uses dynamic programming approach and the second heuristic runs actual influence maximization process while reading influence from the tree.
	
Influence maximization in complex networks through optimal percolation: supplementary information \cite{morone2015}: based on collective theory of optimal influence, this study introduce CI (Collective Influence) algorithm, which identifies maximum influence spreaders by taking into account the weak nodes with strong ties, while also being scalable on large networks. This method also outperforms Degree, Adaptive Degree, PageRank, K-core, Eigenvector and other centrality based heuristics.
	
Influence maximization with viral product design \cite{barbieri2014}: flipping the conventional way of influence maximization by exploiting monotone and sub-modular property of influence mining, this paper  fixes the users and attempts to learn the optimal features, which maximize the spread of a specific product.
	
\todo{Write summary}
IRIE: Scalable and robust influence maximization in social networks \cite{jung2012}: 

\todo{Write summary}
Influence Diffusion Dynamics and Influence Maximization in Social Networks with Friend and Foe Relationships \cite{li2013}: 

\todo{Write summary}
Threshold Models for Competitive Influence in Social Networks \cite{borodin2010}:

\todo{Write summary}
CELF ++: Optimizing the Greedy Algorithm for Influence Maximization in Social Networks \cite{goyal2011b}:

\todo{Write summary}
Social Influence Analysis in Large-scale Networks \cite{tang2009}:


\section*{Learning Influence Probabilities}

\todo{Write summary}
Prediction of Information Diffusion Probabilities for Independent Cascade Model \cite{saito2008}:

\todo{Write summary}
Learning Influence Probabilities in Social Networks \cite{goyal2010}:

\todo{Write summary}
Threshold models of collective behaviour \cite{granovetter1978}:

\section*{Viral Marketing}

\todo{Write summary}
Viral marketing for the real world \cite{watts2007b}:

\todo{Write summary}
Cost-effective Outbreak Detection in Networks \cite{leskovec2007b}:

\todo{Write summary}
Talk of the network: A complex systems look at the underlying process of word-of-mouth \cite{goldenberg2001}:

\todo{Write summary}
Mining the Network Value of Customers \cite{domingos2001}:

\todo{Write summary}
Scalable Influence Maximization for Prevalent Viral Marketing in Large-Scale Social Networks \cite{chen2010b}:

\todo{Write summary}
The Meme Ranking Problem: Maximizing Microblogging Virality \cite{ienco2010}:

\todo{Write summary}
Influence maximization with viral product design \cite{barbieri2014}:

\todo{Write summary}
Mining knowledge-sharing sites for viral marketing \cite{richardson2002}:

\todo{Write summary}
The Dynamics of Viral Marketing \cite{leskovec2007} (Leskovec, Adamic and Huberman, 2007):

\todo{Write summary}
New Product Diffusion Models in Marketing: A Review and Directions for Research (Mahajan, Mullar and Bass, 1990):

\todo{Write summary}
Network-Based Marketing: Identifying Likely Adopters via Consumer Networks (Hill, Provost and Volinsky, 2006):

\todo{Write summary}
Viral Marketing and the Diffusion of Trends on Social Networks (Wortman and Wortman, 2008):

\todo{Write summary}
What exactly is viral marketing (Jurvetson, 2000):

\todo{Write summary}
Conditions for viral influence spreading through multiplex correlated social networks (Hu, Havlin and Makse, 2014):

\todo{Write summary}
Mining Social Networks for Viral Marketing (Domingos, 2005).


\section*{Correlation}

\todo{Write summary}
Influence and Correlation in Social Networks (Anagnostopoulos, Kumar and Mahdian, 2008):

\todo{Write summary}
The tipping point: How little things can make a big difference (Gladwell, 2006):

\todo{Write summary}
Birds of a feather: Homophily in Social Networks (McPherson, Smith-Lovin and Cook, 2001):

\todo{Write summary}
Threshold models of collective behavior (Granovetter, 1978):

\todo{Write summary}
Yes, There is a Correlation - From Social Networks to Personal Behavior on the Web (Singla and Richardson, 2008):

\todo{Write summary}
Randomization tests for distinguishing social influence and homophily effects (La Fond and Neville, 2010):

\todo{Write summary}
Distinguishing influence-based contagion from homophily-driven diffusion in dynamic networks (Aral, Muchnik and Sundararajan, 2009):

\todo{Write summary}
Modeling Relationship Strength in Online Social Networks (Xiang, Neville and Rogati, 2010):

\todo{Write summary}
Feedback Effects between Similarity and Social Influence in Online Communities (Crandall et al., 2008).


\section*{Information Diffusion}

\todo{Write summary}
Influence and Correlation in Social Networks (Anagnostopoulos, Kumar and Mahdian, 2008):

\todo{Write summary}
Back to the Past: Source Identification in Diffusion Networks from Partially Observed Cascades (Farajtabar et al., 2015):

\todo{Write summary}
The Meme Ranking Problem: Maximizing Microblogging Virality (Ienco, Bonchi and Castillo, 2010):

\todo{Write summary}
Influence Diffusion Dynamics and Influence Maximization in Social Networks with Friend and Foe Relationships (Li et al., 2013):

\todo{Write summary}
Viral Marketing and the Diffusion of Trends on Social Networks (Wortman and Wortman, 2008):

\todo{Write summary}
Information diffusion through blogspace (Gruhl et al., 2004):

\todo{Write summary}
Network models of the diffusion of innovations (Valente, 1995):

\todo{Write summary}
Influentials, Networks and Public Opinion Formation (Watts and Dodds, 2007):

\todo{Write summary}
Information diffusion and external influence in networks (Myers, Zhu and Leskovec, 2012):

\todo{Write summary}
Can cascades be predicted? (Cheng et al., 2014):

\todo{Write summary}
Uncovering the Temporal Dynamics of Diffusion Networks (Rodriguez, Balduzzi and Schölkopf, 2011):

\todo{Write summary}
Inferring networks of diffusion and influence (Rodriguez, Leskovec and Krause, 2010):

\todo{Write summary}
A Measurement-driven Analysis of Information Propagation in the Flickr Social Network (Cha, Mislove and Gummadi, 2009):

\todo{Write summary}
Differences in the mechanics of information diffusion across topics: idioms, political hashtags, and complex contagion on twitter (Romero, Meeder and Kleinberg, 2011):

\todo{Write summary}
Personalized recommendation driven by information flow (Song et al., 2006):

\todo{Write summary}
Information flow modeling based on diffusion rate for prediction and ranking (Song et al., 2007):

\todo{Write summary}
Patterns of influence in a recommendation network (Leskovec, Singh and Kleinberg, 2006):

\todo{Write summary}
Absence of influential spreaders in rumor dynamics (Borge-Holthoefer and Moreno, 2012):

\todo{Write summary}
Everyone’s an Influencer: Quantifying Influence on Twitter (Bakshy, Eytan; Mason, Winter A.; Hofman, Jake M.; Watts, 2011).


\section*{Influence Mining Heuristics}

\todo{Write summary}
Centrality in Social Networks, Conceptual Clarification (Freeman, 1979):

\todo{Write summary}
On minimizing budget and time in influence propagation over social networks (Goyal et al., 2012):

\todo{Write summary}
Absence of influential spreaders in rumor dynamics (Borge-Holthoefer and Moreno, 2012):

\todo{Write summary}
Improving the accuracy of the k-shell method by removing redundant links-from a perspective of spreading dynamics (Liu et al., 2015):

\todo{Write summary}
Influential nodes in a diffusion model for social networks (Kempe, Kleinberg and Tardos, 2005):

\todo{Write summary}
Identification of influential spreaders in complex networks (Kitsak et al., 2010):

\todo{Write summary}
Influence Maximization in Social Networks When Negative Opinions May Emerge and Propagate (Chen et al., 2011):

\todo{Write summary}
Identifying Social Influence in Networks Using Randomized Experiments (Aral and Walker, 2011):


\section*{Veracity in Information Diffusion}

\todo{Write summary}
Talk of the network: A complex systems look at the underlying process of word-of-mouth (Goldenberg, Libai and Muller, 2001):

\todo{Write summary}
Back to the Past: Source Identification in Diffusion Networks from Partially Observed Cascades (Farajtabar et al., 2015):

\todo{Write summary}
Trust inference in web-based social networks using resistive networks (Taherian, Amini and Jalili, 2008):

\todo{Write summary}
Information credibility on twitter (Castillo, Mendoza and Poblete, 2011):

\todo{Write summary}
Propagation of trust and distrust (Guha et al., 2004):

\todo{Write summary}
Inferring binary trust relationships in Web-based social networks (Golbeck and Hendler, 2006):

\todo{Write summary}
The Spread of Evidence-Poor Medicine via Flawed Social-Network Analysis (Lyons, 2011):

\todo{Write summary}
Propagation Models for Trust and Distrust in Social Networks (Ziegler and Lausen, 2005):


\section*{Context Awareness}

\todo{Write summary}
Mining Topic-level Influence in Heterogeneous Networks (Liu et al., 2010):

\todo{Write summary}
Maximizing Influence in a Competitive Social Network: A Follower’s Perspective (Carnes et al., 2007):

\todo{Write summary}
Social Influence Analysis in Large-scale Networks (Tang et al., 2009):

\todo{Write summary}
Twitterrank: Finding topic-sensitive influential Twitterers (Weng et al., 2010):

\todo{Write summary}
Mining topic-level opinion influence in microblog (Li et al., 2012):

\todo{Write summary}
Who to Follow and Why: Link Prediction with Explanations (Barbieri, Bonchi and Manco, 2014):

\todo{Write summary}
Topic-Sensitive PageRank (Haveliwala, 2002):

\todo{Write summary}
Content-Based Social Network Analysis of Mailing Lists (Bohn et al., 2011):

\todo{Write summary}
Social Influence and the Diffusion of User-Created Content (Bakshy, Eytan; Karrer, Brian; Adamic, 2009).

\section*{Community Detection}

\todo{Write summary}
CSI: Community-Level Social Influence Analysis (Mehmood and Barbieri, 2013):

\todo{Write summary}
GuruMine: a Pattern Mining System for Discovering Leaders and Tribes (Goyal et al., 2009):

\todo{Write summary}
Learning to discover social circles in ego networks (Leskovec and Mcauley, 2012):

\todo{Write summary}
Discovering leaders from community actions (Goyal, Bonchi and Lakshmanan, 2008):

\todo{Write summary}
Model for generating artificial social networks having community structures with small-world and scale-free properties (Sallaberry, Zaidi and Melançon, 2013):

\todo{Write-summary}
A generalized linear threshold model for an improved description of the spreading dynamics (Ran, Yijun, 2020)

\todo{Write-summary}
Influence maximization: Near-optimal time complexity meets practical efficiency (Youze Tang, Xiaokui Xiao and Yanchen Shi, 2014)

\todo{Write-summary}
Fast and simple decycling and dismantling of networks (Lenka Zdeborová and Pan Zhang, 2016)

